\DocumentMetadata{
  pdfversion=2.0,
  pdfstandard=ua-2,
  lang=en-US,
  tagging=on,
  tagging-setup={math/setup=mathml-SE,tabsorder=structure}
}
\documentclass[11pt,letterpaper]{article}
\usepackage[margin=0.68in,includefoot]{geometry}
\usepackage{newcomputermodern}
\usepackage{amsmath,amscd,mathtools}
\usepackage{tikz,adjustbox}
\providecommand{\square}{\mdlgwhtsquare}
\usepackage[hidelinks]{hyperref}
\hypersetup{
  pdftitle={Variational Analysis/Numerical Optimization PhD Exam, January 2018},
  pdfauthor={University of Florida Department of Mathematics},
  pdfsubject={Graduate mathematics examination},
  pdfdisplaydoctitle=true,
  bookmarksopen=true
}
\setcounter{secnumdepth}{0}
\setlength{\parindent}{0pt}
\setlength{\parskip}{0.28em}
\setlength{\emergencystretch}{3em}
\setlength{\leftmargini}{2.15em}
\setlength{\leftmarginii}{2.65em}
\setlength{\leftmarginiii}{3em}
\renewcommand{\labelenumi}{\textbf{\theenumi.}}
\pagestyle{empty}
\raggedbottom
\allowdisplaybreaks
\newenvironment{examproblems}[1]{%
  \begin{enumerate}%
  \setcounter{enumi}{#1}%
  \setlength{\topsep}{0.35em}%
  \setlength{\partopsep}{0pt}%
  \setlength{\itemsep}{0.48em}%
  \setlength{\parsep}{0.18em}%
}{\end{enumerate}}
\newenvironment{examsubparts}{%
  \begin{enumerate}%
  \setlength{\topsep}{0.18em}%
  \setlength{\partopsep}{0pt}%
  \setlength{\itemsep}{0.16em}%
  \setlength{\parsep}{0.1em}%
}{\end{enumerate}}
\newenvironment{examinstructions}{%
  \par\begingroup\itshape
}{%
  \par\endgroup\vspace{0.18em}
}
\makeatletter
\renewcommand\section{\@startsection{section}{1}{0pt}%
  {-1.25ex plus -.25ex minus -.1ex}%
  {0.55ex plus .1ex}%
  {\normalfont\large\bfseries}}
\renewcommand{\@maketitle}{%
  \newpage\null\vskip -1em%
  {\footnotesize\bfseries
    \href{https://www.ufl.edu/}{University of Florida}, \href{https://math.ufl.edu/}{Department of Mathematics}\par}%
  \vskip -0.5em%
  \tagstructbegin{tag=Title}%
    {\LARGE\bfseries\href{https://gma.math.ufl.edu/exams/numerical-optimization-phd/}{Variational Analysis/Numerical Optimization PhD Exam}, January 2018\par}%
  \tagstructend%
  \vskip 0.25em%
  \tagmcbegin{artifact}%
    \hrule height 1pt%
  \tagmcend%
  \vskip 1em%
}
\makeatother
\title{Variational Analysis/Numerical Optimization PhD Exam, January 2018}
\author{}
\date{}
\begin{document}
\maketitle
\thispagestyle{empty}
\begin{examinstructions}
Do 8 (eight) problems.
\end{examinstructions}
\begin{examproblems}{0}
\item
Let \(\mathbf Q\) denote an~\(n\) by~\(n\) symmetric matrix and let \(\mathbf d_k\), \(1\leq k\leq n\), denote a collection of directions that are \(\mathbf Q\)-conjugate.
\begin{examsubparts}
\item[\textnormal{(a)}]
Suppose the \(\mathbf Q\) is positive definite. Show that the directions \(\mathbf d_k\), \(1\leq k\leq n\), are linearly independent, and express the solution of the linear system \(\mathbf Q\mathbf x=\mathbf b\) in terms of the conjugate directions.
\item[\textnormal{(b)}]
If \(\mathbf d_k^{\mathrm T}\mathbf Q\mathbf d_k>0\) for each~\(k\), then show that \(\mathbf Q\) is positive definite.
\end{examsubparts}
\item
Consider the following variational problem: 
\[
\min\left\{\int_0^1 \frac{1}{2}u'(x)^2+e^{u(x)}\,dx:u\in H^1([0,1]),\quad u(0)=u(1)=0\right\}.
\]
\begin{examsubparts}
\item[\textnormal{(a)}]
What is the first-order necessary optimality condition (Euler equation) for this variational problem?
\item[\textnormal{(b)}]
Consider a uniform mesh \(x_k=kh\) where \(h=1/N\); let \(u_k\) denote an approximation to \(u(x_k)\). Of course, \(u_0=u_N=0\). Give a finite difference approximation in terms \(u_1,\ldots,u_{N-1}\) to the solution of the Euler equation.
\item[\textnormal{(c)}]
Let \(\mathbf F(\mathbf u)\) denote the finite difference system where \(\mathbf u=(u_1,u_2,\ldots,u_{N-1})\in\mathbb R^{N-1}\), and let \(\mathbf u^*\) denote the vector formed by evaluating the solution of the Euler equation at the mesh points \(x_1,x_2,\ldots,x_{N-1}\). Obtain a bound for the components of \(\mathbf F(\mathbf u^*)\) in terms of the mesh spacing~\(h\).
\end{examsubparts}
\item
Consider the following variational problem where~\(f\) is a given smooth function: 
\[
\min\left\{\Phi(u):=\int_0^1 \frac{1}{2}u'(x)^2+f(x)u(x)\,dx:u\in H^1([0,1]),\quad u(0)=u(1)=0\right\}.
\]
 Let \(S^h\) denote the piecewise linear finite element space defined on a uniform mesh with \(v^h(0)=v^h(1)=0\) for all \(v^h\in S^h\). Obtain a bound for the error in the finite element approximation \(u^h\) obtained by minimizing~\(\Phi\) over \(S^h\).
\item
Suppose that \(f:\mathbb R^n\to\mathbb R\) is twice continuously differentiable, \(\mathbf x^*\in\mathbb R^n\), \(\nabla f(\mathbf x^*)=0\), and \(\nabla^2f(\mathbf x^*)\) is positive definite. Show that \(\mathbf x^*\) is a strict local minimizer of~\(f\).
\item
Suppose that a quadratic \(\Phi(\mathbf x)=\frac{1}{2}\mathbf x^{\mathrm T}\mathbf Q\mathbf x-\mathbf b^{\mathrm T}\mathbf x\) is minimized by steepest descent: \(\mathbf x_{k+1}=\mathbf x_k-s_k\mathbf g_k\), where \(\mathbf g_k=\nabla\Phi(\mathbf x_k)\) and \(s_k\) is the stepsize.
\begin{examsubparts}
\item[\textnormal{(a)}]
Derive the formula for the Cauchy step.
\item[\textnormal{(b)}]
Derive the formula for the BB step.
\end{examsubparts}
\item
Consider the system of differential equations \(\dot{\mathbf x}(t)=\mathbf f(\mathbf x(t),\mathbf u(t))\), \(\mathbf x(0)=\mathbf x_0\), where \(\mathbf f:\mathbb R^n\times\mathbb R^m\to\mathbb R^n\) and \(\mathbf f\) satisfies the global Lipschitz condition 
\[
\|\mathbf f(\mathbf x_1,\mathbf u_1)-\mathbf f(\mathbf x_2,\mathbf u_2)\|\leq L\bigl(\|\mathbf x_1-\mathbf x_2\|+\|\mathbf u_1-\mathbf u_2\|\bigr).
\]
 Show that the differential equation has a global solution and on any interval \([0,T]\), the following bound holds for all \((\mathbf u_1,\mathbf u_2)\in\mathcal L^\infty([0,T])\): 
\[
\|\mathbf x_1-\mathbf x_2\|_{\mathcal L^\infty}\leq e^{LT}\|\mathbf u_1-\mathbf u_2\|_{\mathcal L^1},
\]
 where \(\mathbf x_i\) is the solution of the differential equation associated with \(\mathbf u_i\), \(i=1,2\).
\item
Suppose that~\(P\) and~\(Q\) are continuous on \([0,1]\) and the following functional is nonnegative over all \(h\in H_0^1([0,1])\): 
\[
\Omega(h)=\int_0^1 \bigl(P(x)h'(x)^2+Q(x)h(x)^2\bigr)\,dx.
\]
 Show that \(P\geq0\) on \([0,1]\).
\item
Suppose that \(u:[0,1]\to\mathbb R\) is twice continuously differentiable, and let \(u^I\) be the continuous, piecewise linear interpolant of~\(u\). Prove the following bound: 
\[
\|u-u^I\|_{\mathcal L^\infty}\leq\frac{h^2}{8}\|u''\|_{\mathcal L^\infty}.
\]
\item
Suppose that \(f:\mathbb R^n\to\mathbb R\) is a convex function.
\begin{examsubparts}
\item[\textnormal{(a)}]
If~\(f\) is continuously differentiable, then show that for all \(\mathbf x\) and \(\mathbf y\in\mathbb R^n\), we have 
\[
f(\mathbf y)\geq f(\mathbf x)+\nabla f(\mathbf x)(\mathbf y-\mathbf x).
\]
\item[\textnormal{(b)}]
If~\(f\) is twice continuously differentiable, then show that the Hessian \(\nabla^2f(\mathbf x)\) is positive semidefinite for all \(\mathbf x\in\mathbb R^n\).
\end{examsubparts}
\item
Suppose that \(\boldsymbol\Phi:\mathbb R^n\to\mathbb R^n\) and \(\mathbf x^*\in\mathbb R^n\) is a fixed point of \(\boldsymbol\Phi\).
\begin{examsubparts}
\item[\textnormal{(a)}]
If \(\boldsymbol\Phi\) is a contraction mapping with contraction constant~\(\lambda\), then show that the iteration \(\mathbf x_{k+1}=\boldsymbol\Phi(\mathbf x_k)\) converges to \(\mathbf x^*\) from any starting guess \(\mathbf x_0\), and 
\[
\|\mathbf x_k-\mathbf x^*\|\leq\lambda^k\|\mathbf x_0-\mathbf x^*\|.
\]
\item[\textnormal{(b)}]
Suppose that \(\boldsymbol\Phi\) is continuously differentiable on a convex set \(\mathcal K\subset\mathbb R^n\). Let~\(\mu\) denote the supremum of the singular values of \(\nabla\boldsymbol\Phi\) over \(\mathcal K\). Show that 
\[
\|\boldsymbol\Phi(\mathbf x)-\boldsymbol\Phi(\mathbf y)\|\leq\mu\|\mathbf x-\mathbf y\|
\]
 for all \(\mathbf x\) and \(\mathbf y\in\mathcal K\).
\end{examsubparts}
\end{examproblems}
\end{document}
